\begin{abcliste}
\abc
\begin{align}
 p &\approx \peF = \frac{\nch}{4\pi\times\xe} = \pe \approx \result{\peP{0}{2}} \\
 E_\mathrm{kin} &= \EeF = \frac{\left(\pe\right)^2}{2\times\ncme} = \Ee \approx \result{\EevP{-3}{2}}
\end{align}
Small, but never zero: an electron in an atom can never be at rest.

\abc
\begin{align}
 p &\approx \ppF = \frac{\nch}{4\pi\times\xp} = \pp \approx \result{\ppP{0}{2}} \\
 E_\mathrm{kin} &= \EpF = \frac{\left(\pp\right)^2}{2\times\ncmp} = \Ep \approx \result{\EpvP{3}{2}}
\end{align}
Nuclear energies are of this size (\si{\kilo\electronvolt} to \si{\mega\electronvolt}): what holds the nucleus together must be far stronger than the electric forces in an atom.

\abc $p \propto 1/\Delta x$ and $E_\mathrm{kin} \propto p^2 \propto 1/(\Delta x)^2$: half the region gives \result{\text{four times the energy}}.
\end{abcliste}
